%==============================================================
% Analyse des Données — Polycopié complet
% Master GAMMA · Master BA · Licence LMAD — Esprit School of Business
% Tradition : école française (Benzécri, Lebart, Saporta, Escofier-Pagès)
%==============================================================
\documentclass[11pt, a4paper, openany]{book}

% --- Packages ------------------------------------------------
\usepackage[utf8]{inputenc}
\usepackage[T1]{fontenc}
\usepackage[french]{babel}
\usepackage{lmodern}
\usepackage[a4paper, margin=2.5cm]{geometry}
\usepackage{amsmath, amssymb, amsthm, mathtools}
\usepackage{dsfont}
\usepackage{graphicx}
\usepackage{xcolor}
\usepackage{tikz}
\usepackage{pgfplots}
\pgfplotsset{compat=1.18}
\usetikzlibrary{positioning, arrows.meta, shapes, calc}
\usepackage{booktabs}
\usepackage{array}
\usepackage{multirow}
\usepackage{caption}
\usepackage{listings}
\usepackage[hidelinks, colorlinks=true, linkcolor=blue!50!black,
            citecolor=green!40!black, urlcolor=blue!60!black]{hyperref}
\usepackage{enumitem}
\usepackage{fancyhdr}

% --- Couleurs ------------------------------------------------
\definecolor{deepblue}{RGB}{20, 50, 110}
\definecolor{deeporange}{RGB}{220, 100, 30}
\definecolor{darkgreen}{RGB}{30, 110, 60}
\definecolor{lightgray}{RGB}{245, 245, 245}

% --- Listings R ----------------------------------------------
\lstdefinelanguage{Renh}{
  morekeywords={if,else,for,while,function,return,in,TRUE,FALSE,NULL,NA,
                NaN,Inf,library,require,source},
  morecomment=[l]{\#},
  morestring=[b]",
  morestring=[b]'
}
\lstdefinestyle{Rstyle}{
  language=Renh,
  backgroundcolor=\color{lightgray},
  basicstyle=\ttfamily\small,
  keywordstyle=\color{deepblue}\bfseries,
  stringstyle=\color{darkgreen},
  commentstyle=\color{gray}\itshape,
  numbers=left,
  numberstyle=\tiny\color{gray},
  numbersep=8pt,
  frame=single,
  rulecolor=\color{gray!30},
  showstringspaces=false,
  breaklines=true,
  tabsize=2,
  inputencoding=utf8,
  extendedchars=true,
  literate={é}{{\'e}}1 {è}{{\`e}}1 {ê}{{\^e}}1 {à}{{\`a}}1
           {â}{{\^a}}1 {ç}{{\c c}}1 {ô}{{\^o}}1 {î}{{\^\i}}1
           {É}{{\'E}}1 {È}{{\`E}}1 {Ê}{{\^E}}1 {À}{{\`A}}1
           {ù}{{\`u}}1 {û}{{\^u}}1 {ï}{{\"\i}}1 {œ}{{\oe}}1,
}
\lstset{style=Rstyle}

% --- Théorèmes -----------------------------------------------
\theoremstyle{plain}
\newtheorem{theorem}{Théorème}[chapter]
\newtheorem{proposition}[theorem]{Proposition}
\newtheorem{lemma}[theorem]{Lemme}
\newtheorem{corollary}[theorem]{Corollaire}
\theoremstyle{definition}
\newtheorem{definition}[theorem]{Définition}
\newtheorem{example}[theorem]{Exemple}
\newtheorem{exercice}{Exercice}[chapter]
\theoremstyle{remark}
\newtheorem{remark}[theorem]{Remarque}

% --- Macros --------------------------------------------------
\newcommand{\R}{\mathbb{R}}
\newcommand{\N}{\mathbb{N}}
\newcommand{\C}{\mathbb{C}}
\newcommand{\E}{\mathbb{E}}
\newcommand{\1}{\mathds{1}}
\DeclareMathOperator{\Tr}{Tr}
\DeclareMathOperator{\Var}{Var}
\DeclareMathOperator{\Cov}{Cov}
\DeclareMathOperator{\Cor}{Cor}
\DeclareMathOperator{\rang}{rang}
\DeclareMathOperator*{\argmin}{arg\,min}
\DeclareMathOperator*{\argmax}{arg\,max}
\newcommand{\X}{\mathbf{X}}
\newcommand{\V}{\mathbf{V}}
\newcommand{\M}{\mathbf{M}}
\newcommand{\B}{\mathbf{B}}
\newcommand{\W}{\mathbf{W}}
\newcommand{\Z}{\mathbf{Z}}
\newcommand{\U}{\mathbf{U}}
\newcommand{\bSigma}{\mathbf{\Sigma}}
\newcommand{\bLambda}{\mathbf{\Lambda}}
\newcommand{\bg}{\mathbf{g}}
\newcommand{\bx}{\mathbf{x}}
\newcommand{\bu}{\mathbf{u}}
\newcommand{\bv}{\mathbf{v}}

% --- En-têtes ------------------------------------------------
\pagestyle{fancy}
\fancyhf{}
\fancyhead[LE,RO]{\thepage}
\fancyhead[LO]{\itshape\nouppercase{\rightmark}}
\fancyhead[RE]{\itshape\nouppercase{\leftmark}}
\renewcommand{\headrulewidth}{0.4pt}
\setlength{\headheight}{14pt}

\title{
  \vspace{2cm}
  {\Huge \textbf{Analyse des Données}}\\[0.6cm]
  {\Large Théorie, méthodes factorielles et applications}\\[0.5cm]
  {\large ACP $\cdot$ AFC $\cdot$ ACM $\cdot$ Classification $\cdot$ AFD}\\[2cm]
}
\author{
  \Large Master GAMMA · Master Business Analytics\\[0.2cm]
  \large Licence Mathématiques Appliquées (LMAD)\\[0.6cm]
  \large Esprit School of Business
}
\date{\today}

%==============================================================
\begin{document}

\frontmatter
\maketitle

%--------------------------------------------------------------
\chapter*{Préface}
\addcontentsline{toc}{chapter}{Préface}

Ce polycopié accompagne le cours d'\textbf{Analyse des Données} dispensé aux
étudiants de Master GAMMA, Master Business Analytics et Licence Mathématiques
Appliquées (LMAD) de l'Esprit School of Business. Il s'inscrit dans la
\textit{tradition de l'école française d'analyse des données}, fondée par
Jean-Paul Benzécri dans les années 1960 et formalisée par les travaux de
Lebart, Morineau, Saporta, Escofier et Pagès.

L'idée directrice du cours est unique :
\begin{quote}
\textit{Toute méthode factorielle se ramène à une décomposition en valeurs
singulières (SVD) généralisée, dont la lecture géométrique repose sur la
notion d'\textbf{inertie} et le théorème de \textbf{Huygens}.}
\end{quote}

\section*{Public visé}
Ce polycopié s'adresse à des étudiants disposant de connaissances solides en :
\begin{itemize}
  \item algèbre linéaire (valeurs propres, diagonalisation, formes quadratiques);
  \item statistique descriptive et inférentielle (variance, covariance, tests);
  \item programmation R (manipulation de données, graphiques).
\end{itemize}
Les étudiants en Mathématiques Appliquées (LMAD) y trouveront les démonstrations
détaillées, les étudiants en Business Analytics les interprétations et cas
d'usage, et tous bénéficieront du fil conducteur géométrique commun.

\section*{Organisation}
Le cours suit huit chapitres :
\begin{enumerate}
  \item Introduction et statistiques descriptives multivariées;
  \item Algèbre linéaire pour l'analyse des données;
  \item Analyse en Composantes Principales (ACP);
  \item Analyse Factorielle des Correspondances (AFC);
  \item Analyse des Correspondances Multiples (ACM);
  \item Classification (CAH, k-means, HCPC);
  \item Analyse Factorielle Discriminante (AFD);
  \item Méthodes avancées (MDS, CCA, AFM, PLS, KPCA).
\end{enumerate}
Chaque chapitre contient des définitions, théorèmes (avec démonstrations
des plus importants), exemples, code R et exercices.

\section*{Travaux pratiques}
Quatre TPs en R accompagnent le cours :
\begin{itemize}
  \item \textbf{TP1 — ACP} : decathlon (FactoMineR) et vins de Bordeaux;
  \item \textbf{TP2 — AFC} : CSP × Diplôme et Région × Transport;
  \item \textbf{TP3 — ACM} : enquête sur les pratiques numériques;
  \item \textbf{TP4 — AFD} : iris et credit scoring.
\end{itemize}
Tous les jeux de données et leurs scripts de génération sont fournis dans
le dépôt GitHub associé.

\vspace{2cm}
\begin{flushright}
\textit{Esprit School of Business --- \today}
\end{flushright}

\tableofcontents

\mainmatter

%==============================================================
\chapter{Introduction et statistiques descriptives multivariées}
%==============================================================

\section{Position de l'analyse des données}

\begin{definition}[Analyse des données]
L'\textit{analyse des données} (\textit{data analysis}) est l'ensemble des
techniques exploratoires \textit{descriptives}, fondées sur la décomposition
matricielle et la géométrie des nuages de points, visant à révéler la
structure d'un tableau de données $\X \in \R^{n \times p}$ \textit{sans
modèle probabiliste}.
\end{definition}

L'analyse des données se distingue de la statistique inférentielle (qui
postule un modèle probabiliste pour estimer des paramètres et tester des
hypothèses) et du \textit{machine learning} supervisé (qui optimise une
fonction de prédiction).

\paragraph{Objectifs.}
\begin{enumerate}
  \item \textbf{Réduction de dimension} : remplacer $p$ variables corrélées
        par $q \ll p$ composantes synthétiques.
  \item \textbf{Visualisation} : projeter les individus et les variables
        sur des plans interprétables.
  \item \textbf{Classification} (typologie) : regrouper les individus en
        classes homogènes.
\end{enumerate}

\section{Tableau de données}

\begin{definition}[Tableau de données]
On note $\X \in \R^{n \times p}$ le tableau où :
\begin{itemize}
  \item la ligne $i$ représente l'\textit{individu} $\mathbf{x}_i = (x_{i1}, \ldots, x_{ip})^\top$;
  \item la colonne $j$ représente la \textit{variable} $\mathbf{x}^j = (x_{1j}, \ldots, x_{nj})^\top$.
\end{itemize}
\end{definition}

Selon le type de variables, on choisit la méthode adaptée :
\begin{itemize}
  \item Quantitatives $\to$ \textbf{ACP} (chapitre 3);
  \item Tableau de contingence $\to$ \textbf{AFC} (chapitre 4);
  \item Qualitatives nominales $\to$ \textbf{ACM} (chapitre 5);
  \item Mixtes $\to$ AFDM, AFM (chapitre 8).
\end{itemize}

\section{Statistiques descriptives univariées et bivariées}

\subsection{Caractéristiques d'une variable quantitative}
Pour la variable $\mathbf{x}^j$ on définit :
\[
\bar{x}^j = \frac{1}{n}\sum_{i=1}^n x_{ij}, \quad
s_j^2 = \frac{1}{n}\sum_{i=1}^n (x_{ij} - \bar{x}^j)^2, \quad
s_j = \sqrt{s_j^2}.
\]

\begin{remark}
Nous utilisons la variance \textit{biaisée} (divisée par $n$) plutôt que
la variance corrigée ($n-1$) car l'analyse géométrique s'appuie sur les
moments empiriques sans inférence.
\end{remark}

\subsection{Covariance et corrélation}

\begin{definition}[Covariance et corrélation empiriques]
\[
\Cov(\mathbf{x}^j, \mathbf{x}^k) = \frac{1}{n}\sum_{i=1}^n (x_{ij} - \bar{x}^j)(x_{ik} - \bar{x}^k),
\quad
\Cor(\mathbf{x}^j, \mathbf{x}^k) = \frac{\Cov(\mathbf{x}^j, \mathbf{x}^k)}{s_j s_k}.
\]
\end{definition}

\begin{proposition}[Propriétés de la corrélation]
Pour toutes variables $\mathbf{x}^j, \mathbf{x}^k$ :
\begin{enumerate}
  \item $\Cor(\mathbf{x}^j, \mathbf{x}^k) \in [-1, 1]$ (Cauchy-Schwarz);
  \item $|\Cor| = 1 \Leftrightarrow \mathbf{x}^k$ est fonction affine de $\mathbf{x}^j$;
  \item $\Cor$ est invariante par changement d'unités linéaire;
  \item $\Cor = 0 \;\not\Rightarrow\;$ indépendance.
\end{enumerate}
\end{proposition}

\subsection{Forme matricielle}
On note $\X_c = \X - \1_n \bar{\mathbf{x}}^\top$ le tableau \textit{centré}. La
matrice de variance-covariance est
\[
\V = \frac{1}{n} \X_c^\top \X_c \;\in\; \R^{p \times p}.
\]
La matrice de corrélation est
\[
\mathbf{R} = \mathbf{D}_{1/s}\, \V\, \mathbf{D}_{1/s}, \quad
\mathbf{D}_{1/s} = \mathrm{diag}(1/s_j).
\]

\section{Centrage, réduction et tableau standardisé}

\begin{definition}[Tableau réduit]
Le tableau \textit{centré-réduit} (ou \textit{standardisé}) est
\[
\Z = \X_c\, \mathbf{D}_{1/s}, \quad z_{ij} = \frac{x_{ij} - \bar{x}^j}{s_j}.
\]
\end{definition}

\begin{theorem}[Effet sur la variance-covariance]
\[
\V(\X_c) = \V(\X), \qquad \V(\Z) = \mathbf{R}(\X).
\]
\end{theorem}

\begin{proof}
Pour $\X_c$, on a $\bar{\mathbf{x}}_c = 0$, donc $\V(\X_c) = \frac{1}{n}\X_c^\top\X_c = \V(\X)$.
Pour $\Z$, par construction $\Var(z_j) = 1$ et $\Cov(z_j, z_k) = \Cor(x_j, x_k)$,
ce qui donne $\V(\Z) = \mathbf{R}(\X)$.
\end{proof}

\section{Distance, métrique et inertie}

\begin{definition}[Métrique]
Une \textit{métrique} sur $\R^p$ est une matrice $\M$ symétrique définie
positive. La distance associée entre $\mathbf{x}_i$ et $\mathbf{x}_{i'}$ est
\[
d_{\M}^2(\mathbf{x}_i, \mathbf{x}_{i'}) = (\mathbf{x}_i - \mathbf{x}_{i'})^\top \M (\mathbf{x}_i - \mathbf{x}_{i'}).
\]
\end{definition}

\paragraph{Cas particuliers.}
\begin{itemize}
  \item $\M = \mathbf{I}_p$ : distance euclidienne classique;
  \item $\M = \mathbf{D}_{1/s^2}$ : distance euclidienne réduite (ACP normée);
  \item $\M = \V^{-1}$ : Mahalanobis (AFD);
  \item $\M = \mathbf{D}_{1/f_{\bullet j}}$ : distance du $\chi^2$ (AFC).
\end{itemize}

\begin{definition}[Inertie totale d'un nuage]
Soit $\mathcal{N} = \{(\mathbf{x}_i, p_i)\}_{i=1}^n$ avec $\sum_i p_i = 1$ et
$\mathbf{g} = \sum_i p_i \mathbf{x}_i$ le centre de gravité. L'\textit{inertie totale}
de $\mathcal{N}$ par rapport à $\mathbf{g}$ dans la métrique $\M$ est
\[
I_{\mathcal{N}} = \sum_{i=1}^n p_i\, d_{\M}^2(\mathbf{x}_i, \mathbf{g}).
\]
\end{definition}

\begin{proposition}[Inertie et trace]
Avec $\V = \sum_i p_i (\mathbf{x}_i - \mathbf{g})(\mathbf{x}_i - \mathbf{g})^\top$,
\[
I_{\mathcal{N}} = \Tr(\V \M).
\]
\end{proposition}

\begin{proof}
Par cyclicité de la trace, $(\mathbf{x}_i - \mathbf{g})^\top \M (\mathbf{x}_i - \mathbf{g}) = \Tr\bigl(\M (\mathbf{x}_i - \mathbf{g})(\mathbf{x}_i - \mathbf{g})^\top\bigr)$.
En sommant pondéré sur $i$, $\sum_i p_i \Tr\bigl(\M (\mathbf{x}_i - \mathbf{g})(\mathbf{x}_i - \mathbf{g})^\top\bigr) = \Tr(\M \V) = \Tr(\V \M)$.
\end{proof}

\section{Théorème de Huygens}

\begin{theorem}[Huygens]
Soit $\mathcal{N}$ un nuage partitionné en $K$ classes
$\mathcal{C}_1, \ldots, \mathcal{C}_K$ de poids $P_k = \sum_{i \in \mathcal{C}_k} p_i$
et de barycentres $\mathbf{g}_k$. Alors
\[
\underbrace{\sum_i p_i\, d^2(\mathbf{x}_i, \mathbf{g})}_{I_{\text{tot}}}
=
\underbrace{\sum_{k=1}^K \sum_{i \in \mathcal{C}_k} p_i\, d^2(\mathbf{x}_i, \mathbf{g}_k)}_{I_{\text{intra}}}
+
\underbrace{\sum_{k=1}^K P_k\, d^2(\mathbf{g}_k, \mathbf{g})}_{I_{\text{inter}}}.
\]
\end{theorem}

\begin{proof}
Pour tout $\mathbf{x}_i \in \mathcal{C}_k$, on écrit $\mathbf{x}_i - \mathbf{g} = (\mathbf{x}_i - \mathbf{g}_k) + (\mathbf{g}_k - \mathbf{g})$.
Le développement de la norme donne
$d^2(\mathbf{x}_i, \mathbf{g}) = d^2(\mathbf{x}_i, \mathbf{g}_k) + d^2(\mathbf{g}_k, \mathbf{g}) + 2\langle \mathbf{x}_i - \mathbf{g}_k, \mathbf{g}_k - \mathbf{g}\rangle_{\M}$.
Par définition du barycentre $\mathbf{g}_k$, le double produit s'annule en sommant sur
$i \in \mathcal{C}_k$. La pondération sur les classes donne le résultat.
\end{proof}

Ce théorème est fondamental : il fonde le critère du clustering
($\min I_{\text{intra}}$) et de la réduction de dimension
($\max I_{\text{projetée}}$) ainsi que celui de l'AFD
($\max I_{\text{inter}} / I_{\text{intra}}$).

%==============================================================
\chapter{Algèbre linéaire pour l'analyse des données}
%==============================================================

\section{Théorème spectral}

\begin{theorem}[Théorème spectral]
Soit $\mathbf{A} \in \R^{p \times p}$ symétrique réelle. Alors :
\begin{enumerate}
  \item toutes les valeurs propres de $\mathbf{A}$ sont réelles;
  \item les sous-espaces propres associés à des valeurs propres distinctes sont orthogonaux;
  \item il existe une base orthonormée de $\R^p$ formée de vecteurs propres de $\mathbf{A}$.
\end{enumerate}
On a la décomposition
\[
\mathbf{A} = \U\, \bLambda\, \U^\top = \sum_{k=1}^p \lambda_k\, \mathbf{u}_k \mathbf{u}_k^\top,
\]
avec $\U$ orthogonale et $\bLambda = \mathrm{diag}(\lambda_1, \ldots, \lambda_p)$.
\end{theorem}

\begin{proof}[Idée de la démonstration]
On procède par récurrence sur $p$. Pour $p=1$ trivial. Pour $p \geq 2$, on
considère $\mathbf{u}_1$ unitaire maximisant $\mathbf{u}^\top \mathbf{A} \mathbf{u}$ : par compacité
de la sphère, ce maximum existe et vaut $\lambda_1 = \mathbf{u}_1^\top \mathbf{A} \mathbf{u}_1$.
La matrice $\mathbf{A}$ étant symétrique, $\mathbf{u}_1$ est vecteur propre associé.
Sur l'orthogonal $\mathbf{u}_1^\perp$ on applique l'hypothèse de récurrence.
Les valeurs propres sont réelles car $\mathbf{A}$ est auto-adjointe sur $\R^p$.
\end{proof}

\section{Quotient de Rayleigh}

\begin{definition}[Quotient de Rayleigh]
Pour $\mathbf{A}$ symétrique et $\mathbf{u} \neq 0$ :
\[
R_{\mathbf{A}}(\mathbf{u}) = \frac{\mathbf{u}^\top \mathbf{A} \mathbf{u}}{\mathbf{u}^\top \mathbf{u}}.
\]
\end{definition}

\begin{theorem}[Courant-Fischer-Weyl]
Si $\lambda_1 \geq \cdots \geq \lambda_p$ sont les valeurs propres de
$\mathbf{A}$ symétrique, alors
\[
\lambda_1 = \max_{\mathbf{u} \neq 0} R_{\mathbf{A}}(\mathbf{u}),\quad
\lambda_p = \min_{\mathbf{u} \neq 0} R_{\mathbf{A}}(\mathbf{u}),\quad
\lambda_k = \max_{\mathbf{u} \perp \mathbf{u}_1, \ldots, \mathbf{u}_{k-1}} R_{\mathbf{A}}(\mathbf{u}).
\]
\end{theorem}

\begin{proof}
On décompose $\mathbf{u} = \sum_j \alpha_j \mathbf{u}_j$ dans la base orthonormée des vecteurs
propres. Alors $\mathbf{u}^\top \mathbf{A} \mathbf{u} = \sum_j \lambda_j \alpha_j^2$ et
$\mathbf{u}^\top \mathbf{u} = \sum_j \alpha_j^2$. Le quotient s'écrit
\[
R_{\mathbf{A}}(\mathbf{u}) = \frac{\sum_j \lambda_j \alpha_j^2}{\sum_j \alpha_j^2}
\in [\lambda_p, \lambda_1].
\]
Le maximum $\lambda_1$ est atteint pour $\mathbf{u} = \mathbf{u}_1$. Pour le $k$-ième, on impose
$\alpha_1 = \cdots = \alpha_{k-1} = 0$ et le maximum vaut $\lambda_k$.
\end{proof}

\section{Décomposition en valeurs singulières (SVD)}

\begin{theorem}[SVD]
Toute matrice $\X \in \R^{n \times p}$ admet une décomposition
\[
\X = \U\, \bSigma\, \mathbf{V}^\top,
\]
où $\U \in \R^{n \times n}$ et $\mathbf{V} \in \R^{p \times p}$ sont orthogonales,
et $\bSigma \in \R^{n \times p}$ rectangulaire diagonale de coefficients
$\sigma_1 \geq \cdots \geq \sigma_r \geq 0$, $r = \rang(\X)$.
\end{theorem}

\begin{proposition}[Lien spectral]
\[
\X^\top \X = \mathbf{V}\, (\bSigma^\top \bSigma)\, \mathbf{V}^\top, \quad
\X \X^\top = \U\, (\bSigma \bSigma^\top)\, \U^\top.
\]
Les valeurs propres communes non nulles sont $\lambda_k = \sigma_k^2$.
\end{proposition}

\section{Théorème d'Eckart-Young}

\begin{theorem}[Eckart-Young, 1936]
Soit $\X = \sum_{k=1}^r \sigma_k \mathbf{u}_k \mathbf{v}_k^\top$ la SVD de $\X$. Pour tout
$q \leq r$, la matrice $\X_q = \sum_{k=1}^q \sigma_k \mathbf{u}_k \mathbf{v}_k^\top$ est
l'unique matrice de rang $\leq q$ minimisant $\|\X - \mathbf{Y}\|_F$ :
\[
\X_q = \argmin_{\rang(\mathbf{Y}) \leq q} \|\X - \mathbf{Y}\|_F^2,
\quad
\|\X - \X_q\|_F^2 = \sum_{k=q+1}^r \sigma_k^2.
\]
\end{theorem}

\begin{proof}[Esquisse]
Soit $\mathbf{Y}$ de rang $\leq q$. Pour toute matrice $\mathbf{C}$ de rang
$q$, $\|\X - \mathbf{C}\|_F^2 = \sum_k \sigma_k^2(\X - \mathbf{C})$. La meilleure
approximation est obtenue en projetant les composantes singulières dominantes :
on garde les $q$ plus grandes $\sigma_k$ et leurs vecteurs associés. La preuve
formelle utilise l'inégalité de Mirsky pour les normes unitairement invariantes.
\end{proof}

Ce théorème fonde toute la réduction de dimension : la \textit{meilleure
approximation de rang $q$} d'un tableau est la projection sur les $q$
premières directions principales.

\section{SVD généralisée et schéma de dualité}

\begin{theorem}[SVD généralisée]
Soient $\M$ une métrique sur $\R^p$ et $\mathbf{D}_p$ une métrique sur $\R^n$
(toutes deux SDP). Alors $\X$ admet une décomposition
\[
\X = \U\, \bSigma\, \mathbf{V}^\top
\]
avec $\U^\top \mathbf{D}_p \U = \mathbf{I}$ et $\mathbf{V}^\top \M \mathbf{V} = \mathbf{I}$.
Les axes $\mathbf{v}_k$ diagonalisent $\X^\top \mathbf{D}_p \X \M$.
\end{theorem}

Ce cadre unifie l'ACP, l'AFC, l'ACM et l'AFD : chaque méthode correspond
à un choix particulier du \textit{triplet} $(\X, \M, \mathbf{D}_p)$.

%==============================================================
\chapter{Analyse en Composantes Principales (ACP)}
%==============================================================

\section{Formulation et solution}

\paragraph{Cadre.} $\X_c \in \R^{n \times p}$ centré, métrique
$\M = \mathbf{I}_p$ (ACP non normée) ou $\M = \mathbf{D}_{1/s^2}$ (ACP normée),
pondération $p_i = 1/n$. On note $\V = \frac{1}{n} \X_c^\top \X_c$.

\begin{theorem}[Premier axe principal]
Le premier axe principal $\mathbf{u}_1$ est la direction unitaire maximisant
l'inertie projetée :
\[
\mathbf{u}_1 = \argmax_{\|\mathbf{u}\|=1} \mathbf{u}^\top \V \mathbf{u}.
\]
La solution est le vecteur propre dominant de $\V$ et le maximum vaut $\lambda_1$.
\end{theorem}

\begin{proof}
Lagrangien : $\mathcal{L}(\mathbf{u}, \mu) = \mathbf{u}^\top \V \mathbf{u} - \mu(\mathbf{u}^\top \mathbf{u} - 1)$.
La condition $\nabla_{\mathbf{u}} \mathcal{L} = 2(\V \mathbf{u} - \mu \mathbf{u}) = 0$ donne
$\V \mathbf{u} = \mu \mathbf{u}$, et $\mathbf{u}^\top \V \mathbf{u} = \mu$. Le maximum est atteint pour
$\mu = \lambda_1$ et $\mathbf{u} = \mathbf{u}_1$ (vecteur propre dominant).
\end{proof}

\begin{definition}[Composantes principales]
Le $k$-ième axe principal est $\mathbf{u}_k$ vecteur propre de $\V$ associé à $\lambda_k$.
La $k$-ième composante principale est $\mathbf{c}_k = \X_c \mathbf{u}_k \in \R^n$.
\end{definition}

\begin{proposition}[Propriétés des composantes]
\begin{enumerate}
  \item $\bar{c}_k = 0$;
  \item $\Var(\mathbf{c}_k) = \lambda_k$;
  \item $\Cov(\mathbf{c}_k, \mathbf{c}_l) = 0$ pour $k \neq l$;
  \item $\sum_k \lambda_k = \Tr(\V)$.
\end{enumerate}
\end{proposition}

\begin{proof}
(1) Le tableau $\X_c$ est centré donc $\X_c \mathbf{u}_k$ est centré.
(2) $\Var(\X_c \mathbf{u}_k) = \mathbf{u}_k^\top \V \mathbf{u}_k = \lambda_k \mathbf{u}_k^\top \mathbf{u}_k = \lambda_k$.
(3) $\Cov(\X_c \mathbf{u}_k, \X_c \mathbf{u}_l) = \mathbf{u}_k^\top \V \mathbf{u}_l = \lambda_l \mathbf{u}_k^\top \mathbf{u}_l = 0$ (orthogonalité).
(4) $\sum_k \lambda_k = \Tr(\bLambda) = \Tr(\V)$ (par invariance cyclique).
\end{proof}

\section{Représentations graphiques et aides à l'interprétation}

\paragraph{Coordonnées des individus :}
$F_k(i) = \mathbf{x}_i^\top \mathbf{u}_k$ ; on les représente sur le plan factoriel $(F_1, F_2)$.

\paragraph{Coordonnées des variables (cercle des corrélations) :}

\begin{theorem}
En ACP normée, $G_k(\mathbf{x}^j) = \sqrt{\lambda_k}\, u_{kj} = \Cor(\mathbf{x}^j, \mathbf{c}_k)$,
et $\sum_k G_k(\mathbf{x}^j)^2 \leq 1$.
\end{theorem}

\begin{definition}[Aides à l'interprétation]
Pour l'individu $i$ sur l'axe $k$ :
\[
\cos^2_k(i) = \frac{F_k(i)^2}{\|\mathbf{x}_i - \mathbf{g}\|^2}, \quad
\text{ctr}_k(i) = \frac{p_i F_k(i)^2}{\lambda_k}.
\]
\end{definition}

\section{Choix du nombre d'axes}

Trois critères classiques :
\begin{enumerate}
  \item \textbf{Coude} dans le scree plot $\lambda_k$ vs $k$;
  \item \textbf{Kaiser} : retenir les $\lambda_k \geq 1$ (ACP normée);
  \item \textbf{KSS} (Karlis-Saporta-Spinaki) : $\lambda_k > 1 + 2\sqrt{(p-1)/(n-1)}$.
\end{enumerate}

\section{Code R}
\begin{lstlisting}
library(FactoMineR); library(factoextra)
data(decathlon)
res <- PCA(decathlon, scale.unit = TRUE,
           quanti.sup = 11:12, quali.sup = 13, graph = FALSE)
fviz_eig(res, addlabels = TRUE)
fviz_pca_ind(res, col.ind = "cos2", repel = TRUE)
fviz_pca_var(res, col.var = "contrib", repel = TRUE)
\end{lstlisting}

\begin{exercice}
Démontrer que l'ACP normée est équivalente à une ACP non normée du tableau
préalablement standardisé. Conclure que diagonaliser $\mathbf{R}$ revient
à diagonaliser $\V(\Z)$.
\end{exercice}

%==============================================================
\chapter{Analyse Factorielle des Correspondances (AFC)}
%==============================================================

\section{Cadre et notations}

Soit $\mathbf{N} = (n_{ij}) \in \N^{I \times J}$ un tableau de contingence.
On note $n = \sum_{i,j} n_{ij}$, $f_{ij} = n_{ij}/n$, et les marges
$f_{i\bullet}, f_{\bullet j}$.

\begin{definition}[Profil ligne, profil colonne]
$\mathbf{L}_i = (f_{ij}/f_{i\bullet})_j$, $\mathbf{C}_j = (f_{ij}/f_{\bullet j})_i$.
\end{definition}

\section{Distance du chi-deux}

\begin{definition}
\[
d_{\chi^2}^2(\mathbf{L}_i, \mathbf{L}_{i'}) = \sum_j \frac{1}{f_{\bullet j}}
\left(\frac{f_{ij}}{f_{i\bullet}} - \frac{f_{i'j}}{f_{i'\bullet}}\right)^2.
\]
\end{definition}

\begin{theorem}[Inertie totale = $\Phi^2$]
\[
I_{\text{tot}} = \sum_i f_{i\bullet}\, d_{\chi^2}^2(\mathbf{L}_i, \mathbf{L}_\bullet)
= \Phi^2 = \chi^2 / n.
\]
\end{theorem}

\begin{proof}
Le profil moyen ligne est $\mathbf{L}_\bullet = (f_{\bullet j})_j$. On a
\[
d_{\chi^2}^2(\mathbf{L}_i, \mathbf{L}_\bullet) = \sum_j \frac{1}{f_{\bullet j}}
\left(\frac{f_{ij}}{f_{i\bullet}} - f_{\bullet j}\right)^2
= \sum_j \frac{(f_{ij} - f_{i\bullet} f_{\bullet j})^2}{f_{i\bullet}^2 f_{\bullet j}}.
\]
En sommant pondéré par $f_{i\bullet}$ et après simplification,
\[
I_{\text{tot}} = \sum_{i,j} \frac{(f_{ij} - f_{i\bullet} f_{\bullet j})^2}{f_{i\bullet} f_{\bullet j}}
= \frac{1}{n}\sum_{i,j} \frac{(n_{ij} - n f_{i\bullet} f_{\bullet j})^2}{n f_{i\bullet} f_{\bullet j}}
= \chi^2 / n.
\]
\end{proof}

\section{Décomposition factorielle}

L'AFC se ramène à la SVD du tableau standardisé $s_{ij} = (f_{ij} - f_{i\bullet} f_{\bullet j})/\sqrt{f_{i\bullet} f_{\bullet j}}$.

\begin{theorem}[Formules de transition]
\[
F_k(i) = \frac{1}{\sqrt{\lambda_k}}\sum_j \frac{f_{ij}}{f_{i\bullet}}\, G_k(j),
\quad
G_k(j) = \frac{1}{\sqrt{\lambda_k}}\sum_i \frac{f_{ij}}{f_{\bullet j}}\, F_k(i).
\]
\end{theorem}

Ces formules sont la \textit{lecture barycentrique} du biplot AFC : la position
d'une ligne est, à un facteur près, le barycentre des positions des colonnes
pondérées par son profil.

\section{Code R}
\begin{lstlisting}
library(FactoMineR); library(factoextra)
tab <- as.matrix(read.csv("csp_diplome.csv", row.names = 1))
chi <- chisq.test(tab); chi$statistic / sum(tab)   # = Phi^2
res <- CA(tab, graph = FALSE)
fviz_ca_biplot(res, repel = TRUE)
\end{lstlisting}

%==============================================================
\chapter{Analyse des Correspondances Multiples (ACM)}
%==============================================================

\section{Tableau disjonctif complet et tableau de Burt}

\begin{definition}[TDC]
Soit $p$ variables qualitatives, $K = \sum_j K_j$ modalités. Le tableau
disjonctif complet $\Z \in \{0,1\}^{n \times K}$ vérifie $z_{im} = 1$ ssi
l'individu $i$ possède la modalité $m$.
\end{definition}

\begin{definition}[Tableau de Burt]
$\mathbf{B} = \Z^\top \Z \in \N^{K \times K}$. Les blocs diagonaux contiennent
les effectifs des modalités, les blocs hors-diagonaux les tables de contingence
deux à deux.
\end{definition}

\section{ACM = AFC du TDC}

\begin{theorem}[Inertie totale en ACM]
$I_{\text{tot}} = K/p - 1$.
\end{theorem}

\begin{proof}
L'AFC du TDC fait apparaître la marge ligne constante ($z_{i\bullet} = p$) et
la marge colonne $z_{\bullet m} = n_m$. En appliquant la formule de l'inertie
totale de l'AFC après simplification, on trouve $I_{\text{tot}} = K/p - 1$.
\end{proof}

\section{Tableau de Burt et correction de Benzécri}

\begin{theorem}[Lien valeurs propres TDC / Burt]
Soient $\lambda_k^Z$ les valeurs propres de l'AFC du TDC et $\lambda_k^B$
celles de l'AFC du Burt. Alors $\lambda_k^B = (\lambda_k^Z)^2$.
\end{theorem}

\paragraph{Correction de Benzécri.} Seuls les axes $\lambda_k > 1/p$ sont
significatifs. La part d'inertie corrigée est
\[
\tau_k^c = \frac{(\lambda_k - 1/p)^2}{\sum_{l : \lambda_l > 1/p} (\lambda_l - 1/p)^2}.
\]

\section{Code R}
\begin{lstlisting}
library(FactoMineR); library(factoextra)
df <- read.csv("enquete_numerique.csv", stringsAsFactors = TRUE)
res <- MCA(df, quali.sup = 8:9, graph = FALSE)
res$eig
fviz_mca_var(res, repel = TRUE)
\end{lstlisting}

%==============================================================
\chapter{Classification}
%==============================================================

\section{Critère unificateur}

Une partition $\mathcal{C} = \{\mathcal{C}_1, \ldots, \mathcal{C}_K\}$ minimise
$I_{\text{intra}}$, ce qui équivaut à maximiser $I_{\text{inter}}$ par
le théorème de Huygens.

\section{Algorithme k-means}

\begin{definition}[Critère k-means]
$W(\mathcal{C}) = \sum_k \sum_{i \in \mathcal{C}_k} \|\mathbf{x}_i - \mathbf{g}_k\|^2 = n\, I_{\text{intra}}$.
\end{definition}

\begin{theorem}[Convergence de Lloyd]
La suite $W(\mathcal{C}^{(t)})$ est décroissante et minorée par 0, donc
converge en un nombre fini d'itérations vers un minimum local de $W$.
\end{theorem}

\begin{proof}
À chaque itération :
- l'étape d'affectation diminue $W$ (chaque $\mathbf{x}_i$ est affecté au centre le plus proche);
- l'étape de mise à jour minimise $\sum_{i \in \mathcal{C}_k} \|\mathbf{x}_i - \mathbf{g}\|^2$ en $\mathbf{g} = \mathbf{g}_k$ (centre de gravité).
$W$ ne peut prendre qu'un nombre fini de valeurs (il y a $K^n$ partitions),
donc la suite stationnaire après un nombre fini d'étapes.
\end{proof}

\section{Classification ascendante hiérarchique (CAH)}

L'algorithme agglomératif fusionne à chaque étape les deux classes les plus
proches selon une stratégie de \textit{linkage} $\delta$.

\begin{definition}[Indice de Ward]
\[
\delta_{\text{Ward}}(A, B) = \frac{|A| |B|}{|A| + |B|}\, d^2(\mathbf{g}_A, \mathbf{g}_B).
\]
\end{definition}

\begin{theorem}[Interprétation comme perte d'inertie]
La fusion de $A$ et $B$ augmente l'inertie intra-classe de
$\delta_{\text{Ward}}(A, B)$.
\end{theorem}

\begin{proof}
Avant fusion : inertie intra de $A$ et $B$ égale à $\sum_{x \in A} \|x - \mathbf{g}_A\|^2 + \sum_{x \in B} \|x - \mathbf{g}_B\|^2$.
Après fusion : inertie intra de $A \cup B$ vaut la somme précédente plus
$\frac{|A| |B|}{|A| + |B|} d^2(\mathbf{g}_A, \mathbf{g}_B)$, par développement direct
(formule de König). D'où le résultat.
\end{proof}

\section{HCPC : couplage ACP/AFC/ACM + classification}

\paragraph{Procédure.}
\begin{enumerate}
  \item Méthode factorielle (ACP, AFC ou ACM);
  \item Sélection de $q$ axes;
  \item CAH-Ward sur les $q$ premières composantes;
  \item Coupure automatique au plus grand saut;
  \item Consolidation par k-means;
  \item Caractérisation des classes (valeur-test).
\end{enumerate}

\paragraph{Code R.}
\begin{lstlisting}
library(FactoMineR); library(factoextra)
data(decathlon)
res.pca <- PCA(decathlon[, 1:10], graph = FALSE)
res.hcpc <- HCPC(res.pca, nb.clust = -1, consol = TRUE, graph = FALSE)
plot(res.hcpc, choice = "tree")
res.hcpc$desc.var
\end{lstlisting}

%==============================================================
\chapter{Analyse Factorielle Discriminante (AFD)}
%==============================================================

\section{Décomposition de la variance}

Soit $Y \in \{1, \ldots, K\}$ une variable qualitative définissant $K$ groupes
$\mathcal{C}_1, \ldots, \mathcal{C}_K$ d'effectifs $n_k$.

\begin{theorem}[Décomposition $\V = \B + \W$]
\[
\B = \frac{1}{n}\sum_k n_k (\mathbf{g}_k - \mathbf{g})(\mathbf{g}_k - \mathbf{g})^\top, \quad
\W = \frac{1}{n}\sum_k \sum_{i \in \mathcal{C}_k} (\mathbf{x}_i - \mathbf{g}_k)(\mathbf{x}_i - \mathbf{g}_k)^\top.
\]
On a $\V = \B + \W$.
\end{theorem}

\begin{proof}
Pour $\mathbf{x}_i \in \mathcal{C}_k$, $\mathbf{x}_i - \mathbf{g} = (\mathbf{x}_i - \mathbf{g}_k) + (\mathbf{g}_k - \mathbf{g})$.
Le produit extérieur $(\mathbf{x}_i - \mathbf{g})(\mathbf{x}_i - \mathbf{g})^\top$ se développe en :
$(\mathbf{x}_i - \mathbf{g}_k)(\mathbf{x}_i - \mathbf{g}_k)^\top + (\mathbf{g}_k - \mathbf{g})(\mathbf{g}_k - \mathbf{g})^\top + 2 \cdot \text{terme croisé}$.
Le terme croisé s'annule par centrage à $\mathbf{g}_k$ en sommant sur $i \in \mathcal{C}_k$,
ce qui donne le résultat.
\end{proof}

\section{Critère de Fisher}

\begin{definition}[Critère de Fisher]
\[
J(\mathbf{u}) = \frac{\mathbf{u}^\top \B \mathbf{u}}{\mathbf{u}^\top \W \mathbf{u}}.
\]
\end{definition}

\begin{theorem}[Solution]
Les axes discriminants maximisant $J$ sont les vecteurs propres généralisés
$\B \mathbf{u} = \mu \W \mathbf{u}$, soit les vecteurs propres de $\W^{-1} \B$. Le maximum
vaut la plus grande valeur propre $\lambda_1$.
\end{theorem}

\begin{proof}
On contraint $\mathbf{u}^\top \W \mathbf{u} = 1$ et on maximise $\mathbf{u}^\top \B \mathbf{u}$ par
Lagrangien : $\mathcal{L} = \mathbf{u}^\top \B \mathbf{u} - \mu(\mathbf{u}^\top \W \mathbf{u} - 1)$.
La condition $\nabla \mathcal{L} = 0$ donne $\B \mathbf{u} = \mu \W \mathbf{u}$. La valeur
de $J$ atteinte vaut $\mu$, donc le maximum est $\lambda_1$ (plus grande
valeur propre de $\W^{-1} \B$).
\end{proof}

\section{Cas $K = 2$ groupes}

\begin{theorem}[Direction de Fisher pour 2 groupes]
Si $K = 2$, l'unique axe discriminant est, à un scalaire près,
\[
\mathbf{u}^* = \W^{-1}(\mathbf{g}_1 - \mathbf{g}_2).
\]
\end{theorem}

\section{Règle de classement (LDA)}

\begin{definition}[Règle de Mahalanobis]
$\hat{k}(\mathbf{x}^*) = \argmin_k (\mathbf{x}^* - \mathbf{g}_k)^\top \W^{-1} (\mathbf{x}^* - \mathbf{g}_k)$.
\end{definition}

\begin{theorem}[Forme linéaire]
La frontière entre deux classes $k$ et $l$ est l'hyperplan affine
$\{\mathbf{x} : 2(\mathbf{g}_k - \mathbf{g}_l)^\top \W^{-1} \mathbf{x} = \mathbf{g}_k^\top \W^{-1} \mathbf{g}_k - \mathbf{g}_l^\top \W^{-1} \mathbf{g}_l\}$.
\end{theorem}

\section{Optimalité bayésienne (cas gaussien)}

Si $(X | Y = k) \sim \mathcal{N}_p(\mu_k, \Sigma)$ avec covariance commune,
la règle de Bayes coïncide avec l'AFD à un terme de probabilité a priori près.

\begin{theorem}[Bayes/Fisher]
Sous gaussianité et homoscédasticité, $\hat{k}_B(\mathbf{x}) = \argmin_k (\mathbf{x} - \mu_k)^\top \Sigma^{-1} (\mathbf{x} - \mu_k) - 2 \ln \pi_k$,
qui se réduit à la règle d'AFD lorsque $\pi_k = 1/K$.
\end{theorem}

\section{Code R}
\begin{lstlisting}
library(MASS); data(iris)
res <- lda(Species ~ ., data = iris)
predict(res)$x      # coordonnees factorielles discriminantes
res.cv <- lda(Species ~ ., data = iris, CV = TRUE)
table(iris$Species, res.cv$class)
\end{lstlisting}

%==============================================================
\chapter{Méthodes avancées}
%==============================================================

\section{Multidimensional Scaling (MDS)}

\begin{theorem}[Torgerson]
Étant donnée une matrice de distances $\mathbf{D} = (d_{ij})$, soit
$\mathbf{B} = -\frac{1}{2} \mathbf{H} \mathbf{D}^{(2)} \mathbf{H}$ avec
$\mathbf{H} = \mathbf{I}_n - \frac{1}{n}\1_n\1_n^\top$. Si $\mathbf{B}$ est
positive, ses vecteurs propres fournissent des coordonnées $\mathbf{Y}$ telles
que $\|\mathbf{y}_i - \mathbf{y}_j\| = d_{ij}$.
\end{theorem}

\section{Analyse canonique des corrélations (CCA)}

Étant donnés deux groupes $\X \in \R^{n \times p}$ et $\mathbf{Y} \in \R^{n \times q}$,
la CCA cherche $\mathbf{a}, \mathbf{b}$ maximisant
\[
\rho = \Cor(\X \mathbf{a}, \mathbf{Y} \mathbf{b}),
\]
solution de $\V_{XX}^{-1} \V_{XY} \V_{YY}^{-1} \V_{YX} \mathbf{a} = \rho^2 \mathbf{a}$.

\section{Analyse Factorielle Multiple (AFM)}

L'AFM analyse simultanément $J$ tableaux $\X^{(j)}$ portant sur les mêmes
individus, en pondérant chaque bloc par $1/\lambda_1^{(j)}$ pour équilibrer
leur influence.

\section{Régression PLS}

Construit des composantes $\mathbf{t}_k = \X \mathbf{w}_k$ maximisant
$\Cov(\X \mathbf{w}, \mathbf{Y})$. Très utile lorsque $p \gg n$.

\section{Variantes non linéaires}

\begin{itemize}
  \item \textbf{Kernel PCA} : ACP dans un espace de Hilbert via un noyau;
  \item \textbf{Isomap} : MDS sur la distance géodésique d'une variété;
  \item \textbf{LLE} : reconstruction locale par voisinage;
  \item \textbf{t-SNE, UMAP} : préservation des voisinages probabilistes.
\end{itemize}

\section{Conclusion générale}

Toutes les méthodes vues s'inscrivent dans le cadre unique de la
\textbf{SVD généralisée} d'un triplet $(\X, \M, \mathbf{D}_p)$. Le choix
de la métrique $\M$ et de la pondération $\mathbf{D}_p$ définit la méthode :
\[
\begin{array}{ll}
\text{ACP} : (\X_c, \mathbf{I}_p \text{ ou } \mathbf{D}_{1/s^2}, \frac{1}{n}\mathbf{I}_n)
& \text{AFC} : (\text{profils}, \mathbf{D}_{1/f_{\bullet j}}, \mathbf{D}_{f_{i\bullet}}) \\
\text{ACM} : (\Z, \frac{n}{p}\mathbf{D}_{1/n_m}, \frac{1}{n}\mathbf{I}_n)
& \text{AFD} : (\X_c, \W^{-1}, \frac{1}{n}\mathbf{I}_n)
\end{array}
\]

Cette unicité algébrique est une force conceptuelle : elle permet de comprendre,
critiquer et étendre toute nouvelle méthode factorielle. Le théorème de Huygens
fournit l'interprétation géométrique : tout est décomposition d'inertie en
parts expliquée et résiduelle.

%==============================================================
\backmatter

\chapter*{Bibliographie}
\addcontentsline{toc}{chapter}{Bibliographie}

\textbf{Ouvrages fondamentaux (école française)}
\begin{enumerate}
  \item Saporta, G. (2011). \textit{Probabilités, analyse des données et statistique} (3e éd.). Technip.
  \item Lebart, L., Morineau, A., \& Piron, M. (2006). \textit{Statistique exploratoire multidimensionnelle} (4e éd.). Dunod.
  \item Escofier, B., \& Pagès, J. (2008). \textit{Analyses factorielles simples et multiples} (4e éd.). Dunod.
  \item Husson, F., Lê, S., \& Pagès, J. (2017). \textit{Analyse de données avec R} (2e éd.). PUR.
  \item Cornillon, P.-A. et al. (2018). \textit{Statistique avec R} (4e éd.). PUR.
\end{enumerate}

\textbf{Auteurs historiques}
\begin{enumerate}
  \setcounter{enumi}{5}
  \item Benzécri, J.-P. (1973). \textit{L'analyse des données} (2 vol.). Dunod.
  \item Cailliez, F., \& Pagès, J.-P. (1976). \textit{Introduction à l'analyse des données}. SMASH.
  \item Bouroche, J.-M., \& Saporta, G. (1980). \textit{L'analyse des données} (\textit{Que sais-je ?}). PUF.
\end{enumerate}

\textbf{Articles séminaux}
\begin{enumerate}
  \setcounter{enumi}{8}
  \item Pearson, K. (1901). On lines and planes of closest fit. \textit{Phil. Mag.}, 2(11).
  \item Hotelling, H. (1933). Analysis of a complex of statistical variables. \textit{J. Educ. Psych.}, 24.
  \item Fisher, R. A. (1936). The use of multiple measurements in taxonomic problems. \textit{Annals of Eugenics}, 7.
  \item Eckart, C., \& Young, G. (1936). The approximation of one matrix by another of lower rank. \textit{Psychometrika}, 1.
  \item Hotelling, H. (1936). Relations between two sets of variates. \textit{Biometrika}, 28.
  \item Ward, J. H. (1963). Hierarchical grouping to optimize an objective function. \textit{JASA}, 58.
  \item Torgerson, W. S. (1952). Multidimensional scaling: I. Theory and method. \textit{Psychometrika}, 17.
  \item Kruskal, J. B. (1964). Multidimensional scaling. \textit{Psychometrika}, 29.
  \item Wold, H. (1975). Soft modelling by latent variables (PLS). \textit{Perspectives in Probability and Statistics}.
  \item Schölkopf, B., Smola, A., \& Müller, K.-R. (1998). Nonlinear component analysis as a kernel eigenvalue problem. \textit{Neural Computation}, 10(5).
  \item Tenenbaum, J. B. et al. (2000). A global geometric framework for nonlinear dimensionality reduction. \textit{Science}, 290.
  \item Roweis, S. T., \& Saul, L. K. (2000). Nonlinear dimensionality reduction by locally linear embedding. \textit{Science}, 290.
  \item Candès, E. J. et al. (2011). Robust principal component analysis. \textit{JACM}, 58(3).
\end{enumerate}

\textbf{Logiciels et packages R}
\begin{enumerate}
  \setcounter{enumi}{21}
  \item Lê, S., Josse, J., \& Husson, F. (2008). FactoMineR: An R package for multivariate analysis. \textit{JSS}, 25(1).
  \item Kassambara, A., \& Mundt, F. (2020). \textit{factoextra}. R package version 1.0.7.
  \item Dray, S., \& Dufour, A.-B. (2007). The ade4 package. \textit{JSS}, 22(4).
  \item Venables, W. N., \& Ripley, B. D. (2002). \textit{Modern Applied Statistics with S} (4e éd.). Springer.
\end{enumerate}

\end{document}
